% Manuscript skeleton for the llm-factor-mining project.
% Build: make (latexmk + pdflatex + bibtex).  Numbers in the results section come
% from tables/*.tex, generated by ../scripts/render_paper_tables.py from
% ../results/synthetic_benchmark/summary.json.  LLM and real-data results are
% pending and appear only as \todo placeholders; no number may be typed by hand.
\documentclass[11pt]{article}

\usepackage[T1]{fontenc}
\usepackage[utf8]{inputenc}
\usepackage{lmodern}
\usepackage[margin=1in]{geometry}
\usepackage{amsmath}
\usepackage{amssymb}
\usepackage{booktabs}
\usepackage{graphicx}
\usepackage{xcolor}
\usepackage[round]{natbib}
\usepackage[hidelinks]{hyperref}

% Visible placeholders for everything that is not yet evidence.
\newcommand{\todo}[1]{\textcolor{red}{\textbf{[TODO:} #1\textbf{]}}}
\newcommand{\code}[1]{\texttt{#1}}
\newcommand{\E}{\mathbb{E}}
\newcommand{\sign}{\operatorname{sign}}
\newcommand{\IC}{\mathrm{IC}}
\newcommand{\ICIR}{\mathrm{ICIR}}

% Generated from results/synthetic_benchmark/summary.json (synthetic data only).
% Generated by scripts/render_paper_tables.py from results/synthetic_benchmark/summary.json.
% Do not edit by hand. SYNTHETIC DATA - NOT EVIDENCE ABOUT REAL MARKETS.
\newcommand{\SynSymbols}{60}
\newcommand{\SynSessions}{750}
\newcommand{\SynSeeds}{3}
\newcommand{\SynNullSeeds}{10}
\newcommand{\SynBudget}{200}
\newcommand{\SynBatch}{20}
\newcommand{\SynTopK}{5}
\newcommand{\SynFdrAlpha}{0.1}
\newcommand{\SynConfirmAlpha}{0.1}
\newcommand{\SynMinIcDates}{100}
\newcommand{\SynMinIcFraction}{0.5}
\newcommand{\SynTestAlpha}{0.05}
\newcommand{\SynMaxCorr}{0.7}
\newcommand{\SynRecoveryThreshold}{0.7}
\newcommand{\SynTrueDiscoveryThreshold}{0.1}
\newcommand{\SynRecallStride}{5}
\newcommand{\SynSnrLow}{0.05}
\newcommand{\SynSnrHigh}{0.15}
\newcommand{\SynOracleLowMin}{0.0522}
\newcommand{\SynOracleLowMax}{0.0554}
\newcommand{\SynOracleHighMin}{0.1333}
\newcommand{\SynOracleHighMax}{0.1445}
\newcommand{\SynPowerLow}{0.75--1.00}
\newcommand{\SynPowerAll}{0.75--1.00}
\newcommand{\SynRuns}{12}
\newcommand{\SynNullRuns}{20}
\newcommand{\SynIncompleteRuns}{0}
\newcommand{\SynLibraryFamilyFound}{0}
\newcommand{\SynLibraryFamilySelected}{0}
\newcommand{\SynDrawnFound}{0}
\newcommand{\SynDrawnSelected}{0}
\newcommand{\SynDrawnName}{drawn\_40}
\newcommand{\SynDrawnExpression}{\texttt{delta(ts\_cov(ts\_argmin(low,20),cs\_demean(open),5),1)}}
\newcommand{\SynNoveltyReversal}{0.00}
\newcommand{\SynNoveltyVolume}{0.00}
\newcommand{\SynNoveltyLibraryFamily}{0.68--0.72}
\newcommand{\SynNoveltyDrawn}{0.94--0.95}
\newcommand{\SynRuntime}{332.8}
\newcommand{\SynTestIcDates}{148}
\newcommand{\SynRealizedSnrLow}{0.048}
\newcommand{\SynRealizedSnrHigh}{0.142}
\newcommand{\SynIterationsLow}{12--15}
\newcommand{\SynIterationsHigh}{25--32}
\newcommand{\SynMaxResidual}{$9.6\times10^{-13}$}
\newcommand{\SynSoftware}{Python 3.11.15, numpy 2.4.6, pandas 2.3.3, scipy 1.17.1, llm\_factor\_mining 0.2.0}
\makeatletter
% \SynMean{<snr>}{<arm>}{<metric>} and \SynMeanSd{...}: planted-market cells across complete runs.
% \SynNullMean{<arm>}{<metric>}, \SynNullMeanSd{...} and \SynNullSelecting{<arm>}: null markets.
\newcommand{\SynMean}[3]{\ifcsname syn@mean@#1@#2@#3\endcsname\csname syn@mean@#1@#2@#3\endcsname\else\PackageError{synthetic}{no synthetic value #1/#2/#3}{}\fi}
\newcommand{\SynMeanSd}[3]{\ifcsname syn@msd@#1@#2@#3\endcsname\csname syn@msd@#1@#2@#3\endcsname\else\PackageError{synthetic}{no synthetic value #1/#2/#3}{}\fi}
\newcommand{\SynNullMean}[2]{\ifcsname syn@nullmean@#1@#2\endcsname\csname syn@nullmean@#1@#2\endcsname\else\PackageError{synthetic}{no null value #1/#2}{}\fi}
\newcommand{\SynNullMeanSd}[2]{\ifcsname syn@nullmsd@#1@#2\endcsname\csname syn@nullmsd@#1@#2\endcsname\else\PackageError{synthetic}{no null value #1/#2}{}\fi}
\newcommand{\SynNullSelecting}[1]{\ifcsname syn@nullsel@#1\endcsname\csname syn@nullsel@#1\endcsname\else\PackageError{synthetic}{no null arm #1}{}\fi}
\makeatother
\expandafter\def\csname syn@mean@0.05@evolutionary@recovery_rate\endcsname{0.08}
\expandafter\def\csname syn@msd@0.05@evolutionary@recovery_rate\endcsname{0.08 $\pm$ 0.14}
\expandafter\def\csname syn@mean@0.05@evolutionary@recall_rate\endcsname{0.42}
\expandafter\def\csname syn@msd@0.05@evolutionary@recall_rate\endcsname{0.42 $\pm$ 0.14}
\expandafter\def\csname syn@mean@0.05@evolutionary@fdp_true\endcsname{0.00}
\expandafter\def\csname syn@msd@0.05@evolutionary@fdp_true\endcsname{0.00}
\expandafter\def\csname syn@mean@0.05@evolutionary@test_nonsignificant_rate\endcsname{0.50}
\expandafter\def\csname syn@msd@0.05@evolutionary@test_nonsignificant_rate\endcsname{0.50}
\expandafter\def\csname syn@mean@0.05@evolutionary@any_selected\endcsname{0.33}
\expandafter\def\csname syn@msd@0.05@evolutionary@any_selected\endcsname{0.33 $\pm$ 0.58}
\expandafter\def\csname syn@mean@0.05@evolutionary@n_selected\endcsname{1.3}
\expandafter\def\csname syn@msd@0.05@evolutionary@n_selected\endcsname{1.3 $\pm$ 2.3}
\expandafter\def\csname syn@mean@0.05@evolutionary@n_screen_survivors\endcsname{46.7}
\expandafter\def\csname syn@msd@0.05@evolutionary@n_screen_survivors\endcsname{46.7 $\pm$ 14.0}
\expandafter\def\csname syn@mean@0.05@evolutionary@n_confirmed\endcsname{3.7}
\expandafter\def\csname syn@msd@0.05@evolutionary@n_confirmed\endcsname{3.7 $\pm$ 6.4}
\expandafter\def\csname syn@mean@0.05@evolutionary@n_behaviour_classes\endcsname{165.0}
\expandafter\def\csname syn@msd@0.05@evolutionary@n_behaviour_classes\endcsname{165.0 $\pm$ 14.4}
\expandafter\def\csname syn@mean@0.05@evolutionary@n_behaviour_duplicates\endcsname{28.0}
\expandafter\def\csname syn@msd@0.05@evolutionary@n_behaviour_duplicates\endcsname{28.0 $\pm$ 14.2}
\expandafter\def\csname syn@mean@0.05@evolutionary@n_degenerate\endcsname{7.0}
\expandafter\def\csname syn@msd@0.05@evolutionary@n_degenerate\endcsname{7.0 $\pm$ 1.0}
\expandafter\def\csname syn@mean@0.05@evolutionary@composite_test_ic\endcsname{0.0224}
\expandafter\def\csname syn@msd@0.05@evolutionary@composite_test_ic\endcsname{0.0224}
\expandafter\def\csname syn@mean@0.05@evolutionary@composite_test_icir\endcsname{0.162}
\expandafter\def\csname syn@msd@0.05@evolutionary@composite_test_icir\endcsname{0.162}
\expandafter\def\csname syn@mean@0.05@evolutionary@composite_test_t\endcsname{1.88}
\expandafter\def\csname syn@msd@0.05@evolutionary@composite_test_t\endcsname{1.88}
\expandafter\def\csname syn@mean@0.05@evolutionary@selected_test_ic_mean\endcsname{0.0222}
\expandafter\def\csname syn@msd@0.05@evolutionary@selected_test_ic_mean\endcsname{0.0222}
\expandafter\def\csname syn@mean@0.05@evolutionary@composite_test_ls_net_sharpe\endcsname{-0.33}
\expandafter\def\csname syn@msd@0.05@evolutionary@composite_test_ls_net_sharpe\endcsname{-0.33}
\expandafter\def\csname syn@mean@0.05@evolutionary@validity_rate\endcsname{1.00}
\expandafter\def\csname syn@msd@0.05@evolutionary@validity_rate\endcsname{1.00 $\pm$ 0.00}
\expandafter\def\csname syn@mean@0.05@evolutionary@behavioural_novelty_mean\endcsname{0.38}
\expandafter\def\csname syn@msd@0.05@evolutionary@behavioural_novelty_mean\endcsname{0.38}
\expandafter\def\csname syn@mean@0.05@evolutionary@structural_novelty_mean\endcsname{0.88}
\expandafter\def\csname syn@msd@0.05@evolutionary@structural_novelty_mean\endcsname{0.88}
\expandafter\def\csname syn@mean@0.05@random@recovery_rate\endcsname{0.00}
\expandafter\def\csname syn@msd@0.05@random@recovery_rate\endcsname{0.00 $\pm$ 0.00}
\expandafter\def\csname syn@mean@0.05@random@recall_rate\endcsname{0.42}
\expandafter\def\csname syn@msd@0.05@random@recall_rate\endcsname{0.42 $\pm$ 0.14}
\expandafter\def\csname syn@mean@0.05@random@fdp_true\endcsname{0.00}
\expandafter\def\csname syn@msd@0.05@random@fdp_true\endcsname{0.00}
\expandafter\def\csname syn@mean@0.05@random@test_nonsignificant_rate\endcsname{0.00}
\expandafter\def\csname syn@msd@0.05@random@test_nonsignificant_rate\endcsname{0.00}
\expandafter\def\csname syn@mean@0.05@random@any_selected\endcsname{0.33}
\expandafter\def\csname syn@msd@0.05@random@any_selected\endcsname{0.33 $\pm$ 0.58}
\expandafter\def\csname syn@mean@0.05@random@n_selected\endcsname{0.3}
\expandafter\def\csname syn@msd@0.05@random@n_selected\endcsname{0.3 $\pm$ 0.6}
\expandafter\def\csname syn@mean@0.05@random@n_screen_survivors\endcsname{11.7}
\expandafter\def\csname syn@msd@0.05@random@n_screen_survivors\endcsname{11.7 $\pm$ 6.4}
\expandafter\def\csname syn@mean@0.05@random@n_confirmed\endcsname{0.3}
\expandafter\def\csname syn@msd@0.05@random@n_confirmed\endcsname{0.3 $\pm$ 0.6}
\expandafter\def\csname syn@mean@0.05@random@n_behaviour_classes\endcsname{162.0}
\expandafter\def\csname syn@msd@0.05@random@n_behaviour_classes\endcsname{162.0 $\pm$ 2.6}
\expandafter\def\csname syn@mean@0.05@random@n_behaviour_duplicates\endcsname{25.3}
\expandafter\def\csname syn@msd@0.05@random@n_behaviour_duplicates\endcsname{25.3 $\pm$ 2.9}
\expandafter\def\csname syn@mean@0.05@random@n_degenerate\endcsname{12.7}
\expandafter\def\csname syn@msd@0.05@random@n_degenerate\endcsname{12.7 $\pm$ 5.0}
\expandafter\def\csname syn@mean@0.05@random@composite_test_ic\endcsname{0.0186}
\expandafter\def\csname syn@msd@0.05@random@composite_test_ic\endcsname{0.0186}
\expandafter\def\csname syn@mean@0.05@random@composite_test_icir\endcsname{0.131}
\expandafter\def\csname syn@msd@0.05@random@composite_test_icir\endcsname{0.131}
\expandafter\def\csname syn@mean@0.05@random@composite_test_t\endcsname{1.76}
\expandafter\def\csname syn@msd@0.05@random@composite_test_t\endcsname{1.76}
\expandafter\def\csname syn@mean@0.05@random@selected_test_ic_mean\endcsname{0.0186}
\expandafter\def\csname syn@msd@0.05@random@selected_test_ic_mean\endcsname{0.0186}
\expandafter\def\csname syn@mean@0.05@random@composite_test_ls_net_sharpe\endcsname{0.52}
\expandafter\def\csname syn@msd@0.05@random@composite_test_ls_net_sharpe\endcsname{0.52}
\expandafter\def\csname syn@mean@0.05@random@validity_rate\endcsname{1.00}
\expandafter\def\csname syn@msd@0.05@random@validity_rate\endcsname{1.00 $\pm$ 0.00}
\expandafter\def\csname syn@mean@0.05@random@behavioural_novelty_mean\endcsname{0.59}
\expandafter\def\csname syn@msd@0.05@random@behavioural_novelty_mean\endcsname{0.59}
\expandafter\def\csname syn@mean@0.05@random@structural_novelty_mean\endcsname{0.86}
\expandafter\def\csname syn@msd@0.05@random@structural_novelty_mean\endcsname{0.86}
\expandafter\def\csname syn@mean@0.15@evolutionary@recovery_rate\endcsname{0.33}
\expandafter\def\csname syn@msd@0.15@evolutionary@recovery_rate\endcsname{0.33 $\pm$ 0.14}
\expandafter\def\csname syn@mean@0.15@evolutionary@recall_rate\endcsname{0.50}
\expandafter\def\csname syn@msd@0.15@evolutionary@recall_rate\endcsname{0.50 $\pm$ 0.00}
\expandafter\def\csname syn@mean@0.15@evolutionary@fdp_true\endcsname{0.00}
\expandafter\def\csname syn@msd@0.15@evolutionary@fdp_true\endcsname{0.00 $\pm$ 0.00}
\expandafter\def\csname syn@mean@0.15@evolutionary@test_nonsignificant_rate\endcsname{0.00}
\expandafter\def\csname syn@msd@0.15@evolutionary@test_nonsignificant_rate\endcsname{0.00 $\pm$ 0.00}
\expandafter\def\csname syn@mean@0.15@evolutionary@any_selected\endcsname{1.00}
\expandafter\def\csname syn@msd@0.15@evolutionary@any_selected\endcsname{1.00 $\pm$ 0.00}
\expandafter\def\csname syn@mean@0.15@evolutionary@n_selected\endcsname{5.0}
\expandafter\def\csname syn@msd@0.15@evolutionary@n_selected\endcsname{5.0 $\pm$ 0.0}
\expandafter\def\csname syn@mean@0.15@evolutionary@n_screen_survivors\endcsname{85.0}
\expandafter\def\csname syn@msd@0.15@evolutionary@n_screen_survivors\endcsname{85.0 $\pm$ 7.5}
\expandafter\def\csname syn@mean@0.15@evolutionary@n_confirmed\endcsname{76.7}
\expandafter\def\csname syn@msd@0.15@evolutionary@n_confirmed\endcsname{76.7 $\pm$ 7.6}
\expandafter\def\csname syn@mean@0.15@evolutionary@n_behaviour_classes\endcsname{153.3}
\expandafter\def\csname syn@msd@0.15@evolutionary@n_behaviour_classes\endcsname{153.3 $\pm$ 19.6}
\expandafter\def\csname syn@mean@0.15@evolutionary@n_behaviour_duplicates\endcsname{41.3}
\expandafter\def\csname syn@msd@0.15@evolutionary@n_behaviour_duplicates\endcsname{41.3 $\pm$ 19.0}
\expandafter\def\csname syn@mean@0.15@evolutionary@n_degenerate\endcsname{5.3}
\expandafter\def\csname syn@msd@0.15@evolutionary@n_degenerate\endcsname{5.3 $\pm$ 0.6}
\expandafter\def\csname syn@mean@0.15@evolutionary@composite_test_ic\endcsname{0.0762}
\expandafter\def\csname syn@msd@0.15@evolutionary@composite_test_ic\endcsname{0.0762 $\pm$ 0.0126}
\expandafter\def\csname syn@mean@0.15@evolutionary@composite_test_icir\endcsname{0.581}
\expandafter\def\csname syn@msd@0.15@evolutionary@composite_test_icir\endcsname{0.581 $\pm$ 0.074}
\expandafter\def\csname syn@mean@0.15@evolutionary@composite_test_t\endcsname{7.18}
\expandafter\def\csname syn@msd@0.15@evolutionary@composite_test_t\endcsname{7.18 $\pm$ 1.22}
\expandafter\def\csname syn@mean@0.15@evolutionary@selected_test_ic_mean\endcsname{0.0486}
\expandafter\def\csname syn@msd@0.15@evolutionary@selected_test_ic_mean\endcsname{0.0486 $\pm$ 0.0025}
\expandafter\def\csname syn@mean@0.15@evolutionary@composite_test_ls_net_sharpe\endcsname{5.47}
\expandafter\def\csname syn@msd@0.15@evolutionary@composite_test_ls_net_sharpe\endcsname{5.47 $\pm$ 0.91}
\expandafter\def\csname syn@mean@0.15@evolutionary@validity_rate\endcsname{1.00}
\expandafter\def\csname syn@msd@0.15@evolutionary@validity_rate\endcsname{1.00 $\pm$ 0.00}
\expandafter\def\csname syn@mean@0.15@evolutionary@behavioural_novelty_mean\endcsname{0.27}
\expandafter\def\csname syn@msd@0.15@evolutionary@behavioural_novelty_mean\endcsname{0.27 $\pm$ 0.07}
\expandafter\def\csname syn@mean@0.15@evolutionary@structural_novelty_mean\endcsname{0.85}
\expandafter\def\csname syn@msd@0.15@evolutionary@structural_novelty_mean\endcsname{0.85 $\pm$ 0.02}
\expandafter\def\csname syn@mean@0.15@random@recovery_rate\endcsname{0.33}
\expandafter\def\csname syn@msd@0.15@random@recovery_rate\endcsname{0.33 $\pm$ 0.14}
\expandafter\def\csname syn@mean@0.15@random@recall_rate\endcsname{0.42}
\expandafter\def\csname syn@msd@0.15@random@recall_rate\endcsname{0.42 $\pm$ 0.14}
\expandafter\def\csname syn@mean@0.15@random@fdp_true\endcsname{0.00}
\expandafter\def\csname syn@msd@0.15@random@fdp_true\endcsname{0.00 $\pm$ 0.00}
\expandafter\def\csname syn@mean@0.15@random@test_nonsignificant_rate\endcsname{0.13}
\expandafter\def\csname syn@msd@0.15@random@test_nonsignificant_rate\endcsname{0.13 $\pm$ 0.12}
\expandafter\def\csname syn@mean@0.15@random@any_selected\endcsname{1.00}
\expandafter\def\csname syn@msd@0.15@random@any_selected\endcsname{1.00 $\pm$ 0.00}
\expandafter\def\csname syn@mean@0.15@random@n_selected\endcsname{5.0}
\expandafter\def\csname syn@msd@0.15@random@n_selected\endcsname{5.0 $\pm$ 0.0}
\expandafter\def\csname syn@mean@0.15@random@n_screen_survivors\endcsname{58.3}
\expandafter\def\csname syn@msd@0.15@random@n_screen_survivors\endcsname{58.3 $\pm$ 4.2}
\expandafter\def\csname syn@mean@0.15@random@n_confirmed\endcsname{48.0}
\expandafter\def\csname syn@msd@0.15@random@n_confirmed\endcsname{48.0 $\pm$ 8.2}
\expandafter\def\csname syn@mean@0.15@random@n_behaviour_classes\endcsname{162.0}
\expandafter\def\csname syn@msd@0.15@random@n_behaviour_classes\endcsname{162.0 $\pm$ 2.6}
\expandafter\def\csname syn@mean@0.15@random@n_behaviour_duplicates\endcsname{25.3}
\expandafter\def\csname syn@msd@0.15@random@n_behaviour_duplicates\endcsname{25.3 $\pm$ 2.9}
\expandafter\def\csname syn@mean@0.15@random@n_degenerate\endcsname{12.7}
\expandafter\def\csname syn@msd@0.15@random@n_degenerate\endcsname{12.7 $\pm$ 5.0}
\expandafter\def\csname syn@mean@0.15@random@composite_test_ic\endcsname{0.0677}
\expandafter\def\csname syn@msd@0.15@random@composite_test_ic\endcsname{0.0677 $\pm$ 0.0040}
\expandafter\def\csname syn@mean@0.15@random@composite_test_icir\endcsname{0.540}
\expandafter\def\csname syn@msd@0.15@random@composite_test_icir\endcsname{0.540 $\pm$ 0.050}
\expandafter\def\csname syn@mean@0.15@random@composite_test_t\endcsname{6.73}
\expandafter\def\csname syn@msd@0.15@random@composite_test_t\endcsname{6.73 $\pm$ 0.85}
\expandafter\def\csname syn@mean@0.15@random@selected_test_ic_mean\endcsname{0.0415}
\expandafter\def\csname syn@msd@0.15@random@selected_test_ic_mean\endcsname{0.0415 $\pm$ 0.0012}
\expandafter\def\csname syn@mean@0.15@random@composite_test_ls_net_sharpe\endcsname{4.23}
\expandafter\def\csname syn@msd@0.15@random@composite_test_ls_net_sharpe\endcsname{4.23 $\pm$ 1.16}
\expandafter\def\csname syn@mean@0.15@random@validity_rate\endcsname{1.00}
\expandafter\def\csname syn@msd@0.15@random@validity_rate\endcsname{1.00 $\pm$ 0.00}
\expandafter\def\csname syn@mean@0.15@random@behavioural_novelty_mean\endcsname{0.35}
\expandafter\def\csname syn@msd@0.15@random@behavioural_novelty_mean\endcsname{0.35 $\pm$ 0.08}
\expandafter\def\csname syn@mean@0.15@random@structural_novelty_mean\endcsname{0.82}
\expandafter\def\csname syn@msd@0.15@random@structural_novelty_mean\endcsname{0.82 $\pm$ 0.05}
\expandafter\def\csname syn@nullmean@evolutionary@recovery_rate\endcsname{0.00}
\expandafter\def\csname syn@nullmsd@evolutionary@recovery_rate\endcsname{0.00 $\pm$ 0.00}
\expandafter\def\csname syn@nullmean@evolutionary@recall_rate\endcsname{0.42}
\expandafter\def\csname syn@nullmsd@evolutionary@recall_rate\endcsname{0.42 $\pm$ 0.12}
\expandafter\def\csname syn@nullmean@evolutionary@fdp_true\endcsname{n/a}
\expandafter\def\csname syn@nullmsd@evolutionary@fdp_true\endcsname{n/a}
\expandafter\def\csname syn@nullmean@evolutionary@test_nonsignificant_rate\endcsname{n/a}
\expandafter\def\csname syn@nullmsd@evolutionary@test_nonsignificant_rate\endcsname{n/a}
\expandafter\def\csname syn@nullmean@evolutionary@any_selected\endcsname{0.00}
\expandafter\def\csname syn@nullmsd@evolutionary@any_selected\endcsname{0.00 $\pm$ 0.00}
\expandafter\def\csname syn@nullmean@evolutionary@n_selected\endcsname{0.0}
\expandafter\def\csname syn@nullmsd@evolutionary@n_selected\endcsname{0.0 $\pm$ 0.0}
\expandafter\def\csname syn@nullmean@evolutionary@n_screen_survivors\endcsname{4.5}
\expandafter\def\csname syn@nullmsd@evolutionary@n_screen_survivors\endcsname{4.5 $\pm$ 10.9}
\expandafter\def\csname syn@nullmean@evolutionary@n_confirmed\endcsname{0.0}
\expandafter\def\csname syn@nullmsd@evolutionary@n_confirmed\endcsname{0.0 $\pm$ 0.0}
\expandafter\def\csname syn@nullmean@evolutionary@n_behaviour_classes\endcsname{159.7}
\expandafter\def\csname syn@nullmsd@evolutionary@n_behaviour_classes\endcsname{159.7 $\pm$ 16.8}
\expandafter\def\csname syn@nullmean@evolutionary@n_behaviour_duplicates\endcsname{30.7}
\expandafter\def\csname syn@nullmsd@evolutionary@n_behaviour_duplicates\endcsname{30.7 $\pm$ 16.3}
\expandafter\def\csname syn@nullmean@evolutionary@n_degenerate\endcsname{9.6}
\expandafter\def\csname syn@nullmsd@evolutionary@n_degenerate\endcsname{9.6 $\pm$ 3.2}
\expandafter\def\csname syn@nullmean@evolutionary@composite_test_ic\endcsname{n/a}
\expandafter\def\csname syn@nullmsd@evolutionary@composite_test_ic\endcsname{n/a}
\expandafter\def\csname syn@nullmean@evolutionary@composite_test_icir\endcsname{n/a}
\expandafter\def\csname syn@nullmsd@evolutionary@composite_test_icir\endcsname{n/a}
\expandafter\def\csname syn@nullmean@evolutionary@composite_test_t\endcsname{n/a}
\expandafter\def\csname syn@nullmsd@evolutionary@composite_test_t\endcsname{n/a}
\expandafter\def\csname syn@nullmean@evolutionary@selected_test_ic_mean\endcsname{n/a}
\expandafter\def\csname syn@nullmsd@evolutionary@selected_test_ic_mean\endcsname{n/a}
\expandafter\def\csname syn@nullmean@evolutionary@composite_test_ls_net_sharpe\endcsname{n/a}
\expandafter\def\csname syn@nullmsd@evolutionary@composite_test_ls_net_sharpe\endcsname{n/a}
\expandafter\def\csname syn@nullmean@evolutionary@validity_rate\endcsname{1.00}
\expandafter\def\csname syn@nullmsd@evolutionary@validity_rate\endcsname{1.00 $\pm$ 0.00}
\expandafter\def\csname syn@nullmean@evolutionary@behavioural_novelty_mean\endcsname{n/a}
\expandafter\def\csname syn@nullmsd@evolutionary@behavioural_novelty_mean\endcsname{n/a}
\expandafter\def\csname syn@nullmean@evolutionary@structural_novelty_mean\endcsname{n/a}
\expandafter\def\csname syn@nullmsd@evolutionary@structural_novelty_mean\endcsname{n/a}
\expandafter\def\csname syn@nullsel@evolutionary\endcsname{0/10}
\expandafter\def\csname syn@nullmean@random@recovery_rate\endcsname{0.00}
\expandafter\def\csname syn@nullmsd@random@recovery_rate\endcsname{0.00 $\pm$ 0.00}
\expandafter\def\csname syn@nullmean@random@recall_rate\endcsname{0.47}
\expandafter\def\csname syn@nullmsd@random@recall_rate\endcsname{0.47 $\pm$ 0.08}
\expandafter\def\csname syn@nullmean@random@fdp_true\endcsname{n/a}
\expandafter\def\csname syn@nullmsd@random@fdp_true\endcsname{n/a}
\expandafter\def\csname syn@nullmean@random@test_nonsignificant_rate\endcsname{n/a}
\expandafter\def\csname syn@nullmsd@random@test_nonsignificant_rate\endcsname{n/a}
\expandafter\def\csname syn@nullmean@random@any_selected\endcsname{0.00}
\expandafter\def\csname syn@nullmsd@random@any_selected\endcsname{0.00 $\pm$ 0.00}
\expandafter\def\csname syn@nullmean@random@n_selected\endcsname{0.0}
\expandafter\def\csname syn@nullmsd@random@n_selected\endcsname{0.0 $\pm$ 0.0}
\expandafter\def\csname syn@nullmean@random@n_screen_survivors\endcsname{0.2}
\expandafter\def\csname syn@nullmsd@random@n_screen_survivors\endcsname{0.2 $\pm$ 0.6}
\expandafter\def\csname syn@nullmean@random@n_confirmed\endcsname{0.0}
\expandafter\def\csname syn@nullmsd@random@n_confirmed\endcsname{0.0 $\pm$ 0.0}
\expandafter\def\csname syn@nullmean@random@n_behaviour_classes\endcsname{163.2}
\expandafter\def\csname syn@nullmsd@random@n_behaviour_classes\endcsname{163.2 $\pm$ 4.0}
\expandafter\def\csname syn@nullmean@random@n_behaviour_duplicates\endcsname{21.7}
\expandafter\def\csname syn@nullmsd@random@n_behaviour_duplicates\endcsname{21.7 $\pm$ 3.5}
\expandafter\def\csname syn@nullmean@random@n_degenerate\endcsname{15.1}
\expandafter\def\csname syn@nullmsd@random@n_degenerate\endcsname{15.1 $\pm$ 4.0}
\expandafter\def\csname syn@nullmean@random@composite_test_ic\endcsname{n/a}
\expandafter\def\csname syn@nullmsd@random@composite_test_ic\endcsname{n/a}
\expandafter\def\csname syn@nullmean@random@composite_test_icir\endcsname{n/a}
\expandafter\def\csname syn@nullmsd@random@composite_test_icir\endcsname{n/a}
\expandafter\def\csname syn@nullmean@random@composite_test_t\endcsname{n/a}
\expandafter\def\csname syn@nullmsd@random@composite_test_t\endcsname{n/a}
\expandafter\def\csname syn@nullmean@random@selected_test_ic_mean\endcsname{n/a}
\expandafter\def\csname syn@nullmsd@random@selected_test_ic_mean\endcsname{n/a}
\expandafter\def\csname syn@nullmean@random@composite_test_ls_net_sharpe\endcsname{n/a}
\expandafter\def\csname syn@nullmsd@random@composite_test_ls_net_sharpe\endcsname{n/a}
\expandafter\def\csname syn@nullmean@random@validity_rate\endcsname{1.00}
\expandafter\def\csname syn@nullmsd@random@validity_rate\endcsname{1.00 $\pm$ 0.00}
\expandafter\def\csname syn@nullmean@random@behavioural_novelty_mean\endcsname{n/a}
\expandafter\def\csname syn@nullmsd@random@behavioural_novelty_mean\endcsname{n/a}
\expandafter\def\csname syn@nullmean@random@structural_novelty_mean\endcsname{n/a}
\expandafter\def\csname syn@nullmsd@random@structural_novelty_mean\endcsname{n/a}
\expandafter\def\csname syn@nullsel@random\endcsname{0/10}


\title{Searching for Formulaic Alphas with Language Models\\
under Multiple-Testing Control and a Sealed Hold-Out}
\author{\todo{author names and affiliations}}
\date{Working draft --- not submitted}

\begin{document}
\maketitle

\begin{abstract}
Large language models (LLMs) can write formulaic trading signals and justify
them in economic language. Two problems make it hard to credit such signals to
search. First, every candidate is a hypothesis test on one non-repeatable
history, so the number of trials must be counted and controlled. Second, a
model trained on the finance literature may remember which signals worked in
which period. We study whether LLM-guided program search discovers
cross-sectional equity factors that survive multiplicity control over
\emph{every} trial it made, and a single sealed out-of-sample evaluation, more
often than random grammar search and genetic programming at an equal trial
budget. We also study how much of any advantage is explained by memorization.
We contribute an auditable protocol: a typed factor language that cannot look
ahead by construction; an interface that exposes only formation-window
statistics to the proposer; a hash-chained trial ledger that counts invalid
and degenerate proposals as trials; a formation screen over behaviourally
distinct hypotheses followed by a confirmatory test on a validation window no
proposer has seen; a commit-then-reveal test hold-out with a study registry;
and knowledge-cutoff and value-based novelty diagnostics. A synthetic
benchmark with planted signals and null markets validates the harness. With
\SynBudget{} trials per run, neither baseline selected any factor in
\SynNullRuns{} null-market runs, although the formation screen alone passed
spurious expressions for genetic programming. On planted markets both
baselines recovered only textbook signals, never a signal drawn at random
from the grammar, and their selections were mostly substantially correlated
with library factors. \todo{LLM results on the synthetic benchmark and
real-market results on US equities; both are pending and no LLM experiment
has been run.}
\end{abstract}

\section{Introduction}
\label{sec:introduction}

A formulaic alpha is a short program that maps past prices and volumes of a
cross-section of assets to a score predicting next-period relative returns.
Public catalogues of such programs~\citep{kakushadze2016alphas} and decades of
anomaly research show both that the space is rich and that it is easy to
overfit. \citet{harvey2016cross} argue that the conventional $t>2$ hurdle is
far too permissive once the number of tested factors is acknowledged, and
backtest overfitting is a well-documented failure mode in quantitative
finance~\citep{bailey2014deflated,bailey2017pbo}. Automated search worsens
this problem, because a search procedure can generate and test candidates
faster than any human analyst. Genetic
programming~\citep{koza1992genetic} and reinforcement
learning~\citep{yu2023alphagen} have both been used to mine formulaic alphas.

Large language models add a new kind of proposer. An LLM can write
syntactically valid expressions, explain them in economic language, and
refine them in response to feedback. LLM-guided program search has produced
verifiable discoveries in mathematics~\citep{romeraparedes2024funsearch}, and
interactive LLM systems for alpha mining have been
proposed~\citep{wang2023alphagpt}. Finance differs from mathematics in two ways
that matter for evaluation.
\begin{enumerate}
  \item \emph{Ground truth.} There is no verifier that certifies a factor. The
    only evidence is statistical, drawn from one non-repeatable history, so a
    method that proposes more candidates is guaranteed to find more spurious
    ones unless every trial is counted.
  \item \emph{Look-ahead through pre-training.} An LLM may have read which
    signals worked in which period. LLMs have been used to predict returns from
    text~\citep{lopezlira2023chatgpt}, and such predictions can suffer from
    look-ahead bias inherited from the training
    data~\citep{glasserman2023lookahead}. An apparent edge may therefore
    reflect memory rather than search.
\end{enumerate}

\paragraph{Research questions.} This paper asks four questions:
\begin{itemize}
  \item \textbf{RQ1.} At an equal trial budget, does LLM-guided search yield
    frozen factor sets that perform better on a sealed test window than random
    grammar sampling and genetic programming, after multiplicity control that
    counts all trials and confirms on data no proposer has seen?
  \item \textbf{RQ2.} How do the methods differ in search efficiency?
  \item \textbf{RQ3.} How much of any LLM advantage is explained by
    rediscovering textbook factors, or by period-specific knowledge absorbed
    during pre-training?
  \item \textbf{RQ4.} Which components of the LLM proposer matter?
\end{itemize}
The hypotheses and decision rules were written before any LLM or real-data
experiment, but after the first baseline runs on the synthetic benchmark: a
review of those runs led us to revise the benchmark's planted signals and the
synthetic hypothesis, and we disclose this. The rules are given in the
project's research plan, whose hash is to be published before the first LLM
call, and are summarized in Section~\ref{sec:setup}.

\paragraph{Approach.} We treat the proposer as an untrusted component and
build the evaluation protocol around it (Section~\ref{sec:method}):
\begin{itemize}
  \item \emph{The language.} Expressions are written in a typed language
    whose operators only look backward, so no valid expression can read
    future data.
  \item \emph{Search.} The search loop evaluates on a price panel truncated
    at the end of the formation window. The proposer receives only
    formation-window statistics, through closed types whose only free-text
    fields hold the proposal's own text and a validator or fixed harness
    message.
  \item \emph{The ledger.} Every processed proposal is written to a
    hash-chained ledger, including those that fail validation. A formation
    screen counts every budgeted trial and merges behaviourally identical
    ones; because adaptive proposers see formation results, the claim rests
    on a second, confirmatory test on the validation window.
  \item \emph{The hold-out.} The selected set is committed by hash and
    revealed on the test window at most once; a project-wide registry fixes
    each real-data study's windows in advance, counts exploratory
    evaluations and caps reveals.
  \item \emph{Contamination.} Diagnostics compare test performance before and
    after the model's documented knowledge cutoff, and measure distance from
    a library of classic factors on factor \emph{values}.
\end{itemize}

\paragraph{Contributions.}
\begin{enumerate}
  \item An auditable protocol and open implementation for LLM-driven factor
    search. It covers the language, the information barrier, the trial
    ledger, the two-stage multiplicity control, the commit-then-reveal
    hold-out and the study registry, and every control is covered by offline
    tests.
  \item A planted-alpha synthetic benchmark with null markets, whose ground
    truth mixes textbook signals, a documented factor family and a signal
    drawn at random by a procedure fixed in code before the draw was run
    (self-attested, not externally registered); only the drawn signal
    speaks to search beyond prior knowledge. Results on it are reported for
    two baselines.
  \item \todo{An equal-budget comparison of LLM, random-grammar and
    genetic-programming search on US equities (pending).}
  \item \todo{Contamination analysis: knowledge-cutoff difference-in-differences,
    novelty stratification and anonymization probes (pending).}
\end{enumerate}

\paragraph{Current evidence.} The only experiment completed so far is the
harness validation on synthetic data (Section~\ref{sec:results-synthetic}).
With \SynBudget{} trials per run, neither baseline selected a factor in any
of \SynNullRuns{} null-market runs, although genetic programming's formation
screen alone passed spurious expressions. On planted markets both baselines
recovered only textbook signals and never found the drawn signal, and
genetic programming was not clearly better than random sampling. This leaves
room to test whether LLM guidance helps. \todo{State the LLM and real-data findings once the
pre-registered experiments have been run; do not edit this paragraph before
then.}

\section{Related Work}
\label{sec:related}

\paragraph{Formulaic alphas and automated search.}
\citet{kakushadze2016alphas} documents 101 formulaic alphas built from a small
vocabulary of time-series and cross-sectional operators over daily price and
volume data. The factor language in this paper uses a vocabulary of that kind,
and our reference library contains three of its formulas (translated into
the language) and further entries written in that spirit. Searching
over expression trees is the classic setting for genetic
programming~\citep{koza1992genetic}, which serves as our main non-LLM
baseline. \citet{yu2023alphagen} use reinforcement learning to generate
collections of formulaic alphas that work well together. \citet{wang2023alphagpt}
describe an interactive system in which an LLM helps a human analyst mine
alphas. In our study all proposers share the language, the validator's
limits, the trial budget and the evaluation, so that differences between
methods can be attributed to the proposer rather than to the evaluation. The
baselines, however, sample from a restricted menu of windows, literals and
depths, to which the LLM is not limited (Section~\ref{sec:limitations}).

\paragraph{Machine learning in empirical asset pricing.}
\citet{gu2020empirical} compare machine-learning methods for predicting the
cross-section of returns and find gains from flexible models. Our object is
different. We search over short, human-readable programs whose number of
trials can be counted exactly, so that selection effects can be controlled
formally.

\paragraph{LLM-guided program search.}
\citet{romeraparedes2024funsearch} pair an LLM with a programmatic evaluator
and obtain verifiable mathematical discoveries. In finance the evaluator is
itself noisy and the data cannot be regenerated. This moves the burden of
proof from the proposer to the protocol: multiple-testing control and
out-of-sample discipline.

\paragraph{Multiple testing and data snooping.}
We control the false discovery rate with the step-up procedure of
\citet{benjamini1995fdr}. The step-up procedure of \citet{benjamini2001fdr}
serves as a robustness check under arbitrary dependence. Selection-adjusted
Sharpe-type statistics follow the deflated Sharpe ratio
of~\citet{bailey2014deflated}. \citet{harvey2016cross} document the scale of
the multiple-testing problem in the factor literature, and
\citet{bailey2017pbo} formalize the probability of backtest overfitting. The
bootstrap Reality Check of \citet{white2000reality} and the superior
predictive ability test of \citet{hansen2005spa} test whether the best of
many strategies beats a benchmark. The step-down procedure of
\citet{romano2005stepwise} identifies which strategies do so while
controlling the family-wise error rate. We list these tests as planned
extensions, because they need the full trial ledger that our protocol
records. Inference on mean information
coefficients uses heteroskedasticity- and autocorrelation-consistent standard
errors~\citep{newey1987hac}.

\paragraph{Language models, finance and look-ahead.}
\citet{lopezlira2023chatgpt} study whether ChatGPT's reading of news headlines
predicts stock returns. \citet{glasserman2023lookahead} assess look-ahead bias
in GPT-based return predictions. Domain-specific models such as
BloombergGPT~\citep{wu2023bloomberggpt} are trained on large corpora of
financial text. This shows that published financial evidence can enter
pre-training. LLMs are also known to produce fluent but unsupported
content~\citep{ji2023hallucination}, so we treat the economic rationales
attached to proposals as model claims, not as evidence. Our contamination
diagnostics address look-ahead in three ways:
\begin{itemize}
  \item anonymized prompts;
  \item a pre- versus post-cutoff comparison against a baseline;
  \item value-based (behavioural) novelty against a library of known factors.
\end{itemize}

\section{Method}
\label{sec:method}

This section describes the protocol as implemented in the accompanying
package, version~0.2.0. Symbols in \code{typewriter} font name the
corresponding code objects.

\subsection{Data and notation}
\label{sec:data}

A \emph{panel} holds daily bars for trading sessions $t=1,\dots,T$ and assets
$i=1,\dots,N$: open $O_{t,i}$, high $H_{t,i}$, low $L_{t,i}$, close $C_{t,i}$
and volume $V_{t,i}$. The panel accepts only bars whose finality is
\code{confirmed}. For US equities these are read from the suite's shared
store through its market-data gateway; the synthetic benchmark generates bars
in the same canonical schema. Every panel carries a SHA-256 content hash over
its dates, symbols, sources and values. When built from the store, it also
carries the store manifest hash captured at read time.

\subsection{The factor language}
\label{sec:dsl}

\paragraph{Syntax.} Expressions are parsed by a hand-written recursive-descent
parser. User text is never evaluated as code. The grammar is
\begin{align*}
\mathit{expr} &::= \mathit{term}\ \big((\texttt{+}\mid\texttt{-})\ \mathit{term}\big)^{*} &
\mathit{term} &::= \mathit{unary}\ \big((\texttt{*}\mid\texttt{/})\ \mathit{unary}\big)^{*}\\
\mathit{unary} &::= \texttt{-}\,\mathit{unary} \mid \mathit{primary} &
\mathit{primary} &::= \mathit{number} \mid \mathit{name} \mid \mathit{name}\texttt{(}\mathit{expr}(\texttt{,}\,\mathit{expr})^{*}\texttt{)} \mid \texttt{(}\mathit{expr}\texttt{)}
\end{align*}
Infix operators desugar to \code{add}, \code{sub}, \code{mul} and
\code{div}. Unary minus desugars to \code{neg}, or to a negative literal when
applied to a number.

\paragraph{Terminals.} There are eight terminals:
\begin{itemize}
  \item the five bar fields \code{open}, \code{high}, \code{low},
    \code{close} and \code{volume};
  \item \code{returns} $=C_{t}/C_{t-1}-1$, undefined when $|C_{t-1}|<10^{-12}$;
  \item \code{vwap\_proxy} $=(H_t+L_t+C_t)/3$, since daily bars carry no true VWAP;
  \item \code{dollar\_volume} $=C_tV_t$.
\end{itemize}
Forward returns are \emph{not} a terminal. They exist only inside the
evaluation metrics, after the execution lag.

\paragraph{Operators.} The language has 27 operators in three classes
(Table~\ref{tab:operators}). Time-series windows $d$ must be integer
literals: $d\ge 1$ for the shifts \code{delay} and \code{delta}, and
$d\ge 2$ for rolling operators. Rolling operators use only full windows, so a
window touching a missing observation is undefined. Cross-sectional operators
read one date at a time. Non-finite intermediate values are set to missing,
and the final factor is masked to cells with a confirmed bar.

\begin{table}[t]
\centering
\small
\caption{Operators of the factor language. $x,y$ are series, $d$ an integer
window literal and $c$ a numeric literal. Signed transforms preserve the sign
of their argument.}
\label{tab:operators}
\begin{tabular}{lp{0.68\textwidth}}
\toprule
Class & Operators and semantics \\
\midrule
Element-wise (10) & \code{add}, \code{sub}, \code{mul}; \code{div} (undefined where $|y|<10^{-12}$); \code{neg}, \code{abs}, \code{sign}; signed \code{log} $=\sign(x)\log(1+|x|)$; signed \code{sqrt}; signed \code{power}$(x,c)=\sign(x)|x|^{c}$ with $c\in[-4,4]$ \\
Shifts (2) & \code{delay}$(x,d)=x_{t-d}$; \code{delta}$(x,d)=x_t-x_{t-d}$ \\
Rolling (12) & \code{ts\_mean}, \code{ts\_sum}, \code{ts\_min}, \code{ts\_max}; \code{ts\_std} (ddof $1$); \code{ts\_rank} (average-tie percentile rank of $x_t$ among the last $d$ values); \code{ts\_zscore}; \code{ts\_argmax}/\code{ts\_argmin} (sessions since the window extremum, ties to the most recent); \code{decay\_linear} (weights $d,d-1,\dots,1$ from today backwards, normalized by $d(d+1)/2$); \code{ts\_corr}, \code{ts\_cov} (on jointly observed values, ddof $1$; degenerate windows undefined) \\
Cross-sectional (3) & \code{cs\_rank} (per-date percentile rank), \code{cs\_zscore} (ddof $1$), \code{cs\_demean} \\
\bottomrule
\end{tabular}
\end{table}

\paragraph{Static validation.} A validator checks the following and returns
every violation as a machine-readable code, which is fed back to proposers:
\begin{itemize}
  \item operator names, arities and argument kinds;
  \item window and parameter ranges;
  \item depth at most 10 (a leaf has depth $1$);
  \item at most 48 nodes;
  \item windows at most 252 sessions;
  \item literals $|c|\le 10^{6}$;
  \item cumulative lookback at most 504 sessions;
  \item at least one data terminal.
\end{itemize}
The \emph{lookback} of an expression is defined recursively. It is $1$ for
\code{returns}, $0$ for the other terminals and for literals, and for an
operator $f$ with window $d$ applied to arguments $x_j$ it is
\[
\lambda\big(f(x_1,\dots;d)\big)=\max_j \lambda(x_j)+\omega_f(d),
\]
where $\omega_f(d)=d-1$ for rolling operators, $\omega_f(d)=d$ for shifts and
$\omega_f=0$ otherwise.

\paragraph{Causality.} Every operator reads the current row and a bounded
number of earlier rows of its inputs, or a single date across assets.
Structural induction therefore gives the following property: for any valid
expression $e$, the value $e_{t,i}$ is a function of bars dated at or before
$t$ only. The test suite checks this property empirically. It perturbs all
bars after a date $\tau$ and verifies that no factor value at or before
$\tau$ changes, for every operator, every terminal, nested composites and the
reference library.

\paragraph{Canonical form, hashing and complexity.} The canonical form sorts
the two series arguments of the commutative operators (\code{add},
\code{mul}, \code{ts\_corr}, \code{ts\_cov}) by their canonical strings, and
writes literals in non-window slots as floats. The structural hash is the
SHA-256 of a namespaced canonical string. It is the deduplication key and
the evaluator's cache key. Syntax alone cannot identify equivalent
hypotheses: $(a+b)+c$ and $a+(b+c)$ hash differently, and so do
\code{volume}, \code{log(volume)} and \code{cs\_rank(volume)}, whose
per-date ranks are identical. Equivalence for multiple testing is therefore
decided on behaviour (Section~\ref{sec:protocol}). The complexity of an
expression is
$\kappa(e)=n_{\mathrm{op}}+n_{\mathrm{term}}+\tfrac12 n_{\mathrm{lit}}$.

\subsection{Evaluation timing and metrics}
\label{sec:metrics}

\paragraph{Execution lag.} A signal $f_{t,i}$ observed at the close of $t$ is
traded at the close of $t+\ell$ and held for $h$ sessions. The forward
return is
\[
y^{(h)}_{t,i}=\frac{C_{t+\ell+h,i}}{C_{t+\ell,i}}-1,\qquad \ell\ge 1,
\]
and $\ell\ge1$ is enforced in code. All metric series are indexed by the
signal date $t$. We use $\ell=h=1$ throughout.

\paragraph{Windows and embargo.} The calendar is split chronologically by
session counts into formation, validation and test windows. The default
fractions are $0.6/0.2/0.2$ of the sessions left after two embargo gaps. The
windows are separated by $E\ge\ell+h_{\max}$ sessions that belong to no
window. Within a window $W$ with last session position $p_W$, only signal
positions with $p+\ell+h\le p_W$ are scored, so no forward return ends
outside its window. Because factors are causal, each phase computes a factor
on the panel it is allowed to see and then restricts the result to the
window's scored dates. The value on any date is the same whichever panel was
used, so there is no recomputation effect at window boundaries.

\paragraph{Information coefficient.} The per-date rank IC is the Pearson
correlation of average-tie ranks of $f_{t,\cdot}$ and $y_{t,\cdot}$ over
assets where both are observed. It is undefined on dates with fewer than
$10$ such assets or no rank variation. For the IC series on a window's
scored dates:
\begin{itemize}
  \item $\overline{\IC}$ is the mean;
  \item $\ICIR=\overline{\IC}/\mathrm{sd}(\IC)$ is per period, not
    annualized;
  \item $t_{\mathrm{NW}}=\overline{\IC}/\widehat{\mathrm{se}}_{\mathrm{NW}}$
    uses the Newey--West standard error with a Bartlett
    kernel~\citep{newey1987hac} and $q=\max\{h-1,\lfloor 4(n/100)^{2/9}\rfloor\}$
    lags, where $n$ is the number of IC dates.
\end{itemize}
p-values refer $t_{\mathrm{NW}}$ to a $t$ distribution with $n-1$ degrees of
freedom, a conservative small-sample convention.

\paragraph{Portfolio metrics.} On each date, assets are split into five
equal-count buckets by signal. The long-short portfolio holds $+1$ spread
equally over the top bucket and $-1$ over the bottom bucket. With $h>1$, the
weights held are the average of the last $h$ formations. Weights held on
signal date $t$ earn the close-to-close return from $t+\ell$ to $t+\ell+1$.
Trading costs are $10$ basis points per unit of $\sum_i|w_{t,i}-w_{t-1,i}|$.
Sharpe ratios are annualized with $\sqrt{252}$. Top-bucket turnover is the
share of today's top-bucket names that were not in yesterday's top bucket.

\subsection{Search, selection and the sealed hold-out}
\label{sec:protocol}

\paragraph{Phase 1: search.} The search loop evaluates proposals on a panel
truncated at the end of the formation window, so every statistic a proposer
can receive is computed from formation data. The harness can also be given a
loader instead of a panel; each phase then reads only its own data, and the
test window is loaded only after the factor set is frozen. The store path of
the command-line interface works this way. Each round proceeds as follows.
\begin{enumerate}
  \item The proposer receives a context (described below) and returns a batch
    of $b=20$ expressions, or the remaining budget if smaller.
  \item Each expression is validated. An invalid expression becomes a
    \emph{trial} with its error codes. If its whitespace- and
    case-normalized text repeats an earlier invalid text, it is logged as a
    duplicate instead.
  \item A valid expression is canonicalized. If its hash was seen before, it
    is logged as a duplicate and consumes no budget. Otherwise it is
    evaluated on the formation window and becomes a trial. A numerical
    failure is also a trial, with status \code{error}.
  \item An evaluated trial whose IC is defined on fewer than
    $\max(\SynMinIcDates, \lceil \SynMinIcFraction\, n_{\mathrm{scorable}}\rceil)$
    formation dates is \emph{degenerate}: it is a trial with $p=1$, and the
    proposer is told so. Without this rule, and with a normal reference, a
    factor with three defined IC dates can obtain an arbitrarily small
    p-value.
\end{enumerate}
The loop stops for any of these reasons:
\begin{itemize}
  \item the budget of $B$ trials is reached;
  \item $\max(100, 10\lceil B/b\rceil)$ rounds have run (100 at $B=200$,
    $b=20$), so that larger budgets can be used in full;
  \item 5 consecutive rounds produced no new trial;
  \item 3 consecutive proposer failures occurred (for example, refusals).
\end{itemize}
Proposer failures are logged and do not count as trials. Only a run that
reaches its budget is \emph{complete}. A run that ends with repeated proposer
failures, or with no trial at all, is \emph{aborted}: it is logged but
neither committed nor revealed. Misconfiguration (a rejected request, key,
permission or model identifier) is not a proposer failure; it stops the run
with an error, so it cannot pass for a null result.

\paragraph{Information barrier.} The barrier is an interface: it limits what
the harness passes to the proposer. It is not a sandbox; in-process code
could read anything in memory. For the LLM arm, only the rendered prompt
leaves the process. The context type is frozen and slotted. Its feedback
fields accept only two exact types, and subclasses are rejected:
\begin{itemize}
  \item evaluated feedback: canonical text, formation IC, ICIR,
    $t_{\mathrm{NW}}$, top-bucket turnover and complexity. The only factory
    that builds it from a metrics report refuses any window not named
    \code{formation};
  \item rejected feedback: the proposal's text, its error codes and a
    message, truncated to 400 characters. For an invalid proposal the message
    is the validator's; a failed or degenerate trial gets a fixed text (for a
    failure, only the exception type) without the date counts that the
    ledger records.
\end{itemize}
Each round's context contains:
\begin{itemize}
  \item the grammar card and the active limits;
  \item the 8 trials with the highest $|\ICIR|$;
  \item the 4 with the lowest, once more than 8 exist;
  \item the previous round's evaluations (the LLM prompt shows at most 10);
  \item the last 8 rejections;
  \item the trial count and the remaining budget.
\end{itemize}
It contains no dates, asset identifiers, date counts, or validation or test
quantities. The formation statistics still imply the approximate formation
sample size, because $t_{\mathrm{NW}}/\ICIR$ grows roughly like $\sqrt{n}$. An end-to-end test perturbs all post-formation bars and confirms
that every context a proposer receives is unchanged.

\paragraph{Phase 2: selection.} Selection loads a panel truncated at the end
of the validation window.
\begin{enumerate}
  \item \emph{Behavioural classes.} Two evaluated trials whose per-date
    cross-sectional ranks on the formation IC dates agree up to sign have
    identical rank-IC series against any forward return, so they test one
    hypothesis. Each such class is represented by its first trial. Counting
    the copies separately would make Benjamini--Hochberg \emph{less} strict,
    because a step-up rule rewards clusters of equal p-values.
  \item \emph{Formation screen.} The Benjamini--Hochberg step-up
    procedure~\citep{benjamini1995fdr} at $q=\SynFdrAlpha$ is applied to the
    two-sided formation p-values of the class representatives, with family
    size $m=\max(B,n_{\mathrm{trials}})-n_{\mathrm{dup}}$: every budgeted
    slot that did not repeat a tested hypothesis counts, degenerate, invalid
    and failed trials and unused budget as $p=1$, so a run that stops early
    faces the family of a full one. A feedback-driven proposer chooses its
    next expressions after seeing formation results, so these p-values are
    not valid for false-discovery control. The screen is a filter only.
    The implementation also offers
    Benjamini--Yekutieli~\citep{benjamini2001fdr} and Holm's step-down
    procedure. \todo{add a citation for Holm's procedure once its metadata
    has been verified.}
  \item \emph{Orientation and confirmation.} Each survivor $j$ is multiplied
    by $s_j=-1$ if its formation IC is negative and $s_j=+1$ otherwise, and
    scored on the validation window, which no proposer has seen.
    Benjamini--Hochberg at $q=\SynConfirmAlpha$ is applied to the one-sided
    validation p-values of all survivors, with $m$ equal to the number of
    survivors; a survivor that is degenerate on validation gets $p=1$. This
    is the step whose error control carries the claim: given the candidate
    set, which depends on formation data only, the validation p-values are
    valid.
  \item \emph{Decorrelation.} Confirmed candidates are taken in descending
    order of validation ICIR, with ties broken by hash. A candidate is kept
    if the mean over validation dates of its cross-sectional Spearman
    correlation with every kept factor is below $0.7$ in absolute value. At
    most $k=5$ factors are kept.
\end{enumerate}
For each survivor we also record the deflated Sharpe
ratio~\citep{bailey2014deflated} of its formation ICIR as a
\emph{diagnostic}; it is not used for selection. Survivors are selected on
$|\ICIR|$, a two-sided choice: under a symmetric null the maximum of $N$
absolute values behaves like the maximum of $2N$ signed ones, so the null
threshold uses $2N$ trials,
\[
\widehat{SR}_0=\sqrt{V}\Big[(1-\gamma)\,\Phi^{-1}\!\big(1-\tfrac1{2N}\big)+\gamma\,\Phi^{-1}\!\big(1-\tfrac{1}{2Ne}\big)\Big],
\]
where $V$ is the sample variance of the class representatives' formation
ICIRs, $N$ the screen's family size and $\gamma$ the Euler--Mascheroni
constant. The diagnostic is then
\[
\mathrm{DSR}=\Phi\!\left(\frac{(\widehat{SR}-\widehat{SR}_0)\sqrt{n-1}}{\sqrt{1-g_3\widehat{SR}+\tfrac{g_4-1}{4}\widehat{SR}^2}}\right),
\]
with $\widehat{SR}=|\ICIR|$, and $g_3$, $g_4$ the skewness (sign-adjusted by
$s_j$) and non-excess kurtosis of the IC series. Treating trials as
independent makes this diagnostic approximate.

\paragraph{Phase 3: commit and reveal.} The selected set $F$ is an ordered
tuple of canonical expressions, hashes and orientations. The seal
identifier is $\sigma=H(\text{terms})$, where $H$ is SHA-256 over canonical
JSON and the terms are the test window, the panel content hash, the metric
configuration and the language limits. The commitment
$c=H(\sigma,\text{terms},F,H(F))$ is appended to the ledger before any test
metric exists. The one-commit, one-reveal rule is keyed on a \emph{guard}
$g=H(\text{data hash},\text{test-window dates})$ only, so changing the metric
settings or the limits yields a new $\sigma$ but not a second reveal of the
same window. The ledger is re-read on every commit and reveal, so several
instances, or several writers of one ledger file, share the rule. A commit is
refused if the guard already holds a commitment or a reveal, or if the data
hash recomputed from the panel's values differs from the one it was built
with. A reveal recomputes the commitment for the set it is given and refuses
in four cases:
\begin{itemize}
  \item nothing has been committed;
  \item a reveal for $g$ already exists in the ledger;
  \item the set or the seal terms differ from the committed ones;
  \item an external authority refuses it (below).
\end{itemize}
The last three refusals are written to the ledger. On a successful reveal,
test metrics are computed for each factor and for the composite. The
composite is the equal-weight, missing-aware mean of per-date z-scored
oriented signals.

\paragraph{Study registry.} Across runs, a project-wide register fixes each
real-data \emph{study} (symbols, dates, split and a pre-registered maximum
number of reveals) before anything is evaluated, counts every exploratory
evaluation outside a search, records commitments and caps reveals. On store
data the search only commits; revealing is a separate command that requires
the commitment hash, which is meant to be published outside the repository
first, and the registry's permission. The registry, like the ledger, is
self-attested: it makes repeated reveals countable, not impossible.

\paragraph{Trial ledger.} The ledger is an append-only JSON Lines file. Record
$k$ stores
\[
h_k=H(\text{version},k,\text{kind}_k,\tau_k,\text{payload}_k,h_{k-1}),\qquad h_{-1}=0^{64},
\]
where $\tau_k$ is an optional timestamp. Canonical JSON uses sorted keys and
rejects NaN. The ledger records the run configuration, every trial,
duplicate, proposer error and round, the selection with all candidates and
decisions, the commitment, the reveal and the run end. Any edit, insertion,
deletion or reordering breaks the chain. Truncating the tail is detected by
comparing against a separately stored head hash. These checks detect
accidental or naive edits only: the head is written by the same process,
into the same directory, so anyone with write access can drop records and
rebuild the chain. A head or commitment hash proves something to a third
party only once it has been published elsewhere before the reveal.

\subsection{Proposers}
\label{sec:proposers}

\paragraph{Random grammar (A).} This baseline samples typed trees i.i.d. and
ignores all feedback.
\begin{itemize}
  \item Maximum depth is 4.
  \item A series slot below the root becomes a leaf with probability $0.35$.
  \item The second argument of an arithmetic operator becomes a literal from
    $\{0.5,1,2\}$ with probability $0.15$.
  \item Windows are drawn from $\{1,2,3,5,10,20,40,60\}$ (shifts may use
    $1$; rolling operators need at least $2$) and exponents from
    $\{0.5,2\}$. The sampling probability of any tree is computable, which
    lets a test confirm that every planted benchmark signal is in the
    sampler's support.
  \item Sampling repeats until the tree is valid. It is then redrawn, up to
    50 times, if it repeats one of the proposer's earlier samples.
\end{itemize}

\paragraph{Genetic programming (B).} This baseline follows the classic
setup~\citep{koza1992genetic}. It keeps an archive of evaluated trees with
fitness $|\ICIR_{\mathrm{formation}}|-0.002\,\kappa$; the sign is free
because selection orients factors. Each round:
\begin{itemize}
  \item the 30 fittest archive members form the population;
  \item children are bred by tournament selection (size 3) and one of three
    operations: subtree crossover ($p=0.5$), subtree mutation ($0.3$, with
    new subtrees of depth at most 3) or point mutation ($0.2$), which swaps
    an operator for another of the same signature, a terminal, or a window;
  \item a tenth of each batch is random immigrants.
\end{itemize}
Children must be valid and new. These operator settings were fixed before
any experiment and were not tuned on results. The window menu that both
baselines share was extended by the value $1$ in the revision disclosed in
Section~\ref{sec:limitations}, after the first harness runs had been seen,
so that the library-family planted signal's one-session change became
reachable; before, it was $\{2,3,5,10,20,40,60\}$.

\paragraph{LLM-guided (D).} The prompts are two versioned templates, a system
prompt and a user prompt (version~2), whose SHA-256 hashes are recorded with
every call and every proposal. Version~2 dropped the example mechanism labels
of version~1, which named the stories behind two planted benchmark signals.
The system prompt states four things:
\begin{itemize}
  \item the data are anonymized;
  \item the model should not try to infer the market or period;
  \item signs are free;
  \item short, economically motivated expressions are preferred.
\end{itemize}
The user prompt carries the grammar card, the formation feedback blocks
(best, worst and previous round), the recent rejections and the remaining
budget. The response must match a JSON schema with exactly the fields
\code{expression}, \code{rationale} and \code{economic\_mechanism} per
proposal, and no additional properties.

Before sending, the harness-written text of every prompt (templates and
grammar card) is checked for:
\begin{itemize}
  \item ISO-style dates;
  \item month names (full names in any case; ``March'', ``May'' and
    abbreviations when capitalized);
  \item four-digit year-like numbers, excluding the integer counters;
  \item any whole-word occurrence of the panel's symbols.
\end{itemize}
A violation aborts the run. Model-written text echoed back in the rejection
block is checked for the same patterns (except years) and redacted rather
than aborting the run, since it is the model's own output.

The Claude backend sends the following:
\begin{itemize}
  \item the model identifier (default \code{claude-opus-5}) and a maximum of
    16{,}000 output tokens;
  \item adaptive thinking;
  \item an effort level (default \code{high}) and the JSON-schema output
    format.
\end{itemize}
It sends no sampling parameters and no prefilled assistant turn. Responses
are handled as follows:
\begin{itemize}
  \item a refusal stop reason is checked before content is read;
  \item rate limits, overload and server errors, and connection failures are
    retried with bounded exponential backoff;
  \item a rejected request, key, permission or model identifier (HTTP 400,
    401, 403, 404, 413 or 422) stops the run;
  \item truncation at the token limit, invalid JSON and failures that
    survived the retries raise recoverable proposer errors;
  \item paused turns are resumed up to three times.
\end{itemize}
The served model identifier, stop reason, token usage and request
identifier are recorded for every call. Server-side model fallback can
silently change the model under test, so it is an explicit option that is
off in all experiments. A recording layer appends every request and its
outcome to a write-only JSON Lines file, in call order and failures included.
It is keyed by the hash of the request, including a replicate tag and a tag
for the benchmark cell, and of the backend identity. A replay backend
answers the $n$-th occurrence of a request with the $n$-th recorded outcome
and treats anything unrecorded as fatal, so a replayed run follows the live
run's path exactly; a test checks this for a live run that met a transient
failure.

\subsection{Contamination diagnostics}
\label{sec:contamination}

\paragraph{Knowledge-cutoff split.} Given a per-date IC series and a cutoff
session $c$ from the model provider's documentation, the split works as
follows:
\begin{itemize}
  \item \emph{pre} uses signal positions with $p+\ell+h<c$;
  \item \emph{post} uses positions with $p\ge c$;
  \item signal dates whose forward return straddles the cutoff are dropped.
\end{itemize}
The difference in mean IC is standardized by
$\sqrt{\widehat{\mathrm{se}}_{\mathrm{pre}}^2+\widehat{\mathrm{se}}_{\mathrm{post}}^2}$,
with Newey--West standard errors. Section~\ref{sec:setup} uses this split in
a difference-in-differences design against a baseline. No cutoff date is
assumed anywhere in the code.

\paragraph{Novelty.} Rediscovery is measured on values. The
\emph{behavioural novelty} of a factor $f$ against the reference library
$\mathcal{L}$ is
\[
1-\max_{\ell\in\mathcal{L}}\Big|\frac{1}{|T|}\sum_{t\in T}\rho_S\big(f_{t,\cdot},\ell_{t,\cdot}\big)\Big|,
\]
the largest mean cross-sectional Spearman correlation with any library
entry over the formation IC dates $T$. A re-spelling of a classic factor,
such as \code{-ts\_sum(returns, 5)} for five-day reversal, scores near
$0$. \emph{Structural novelty}, reported as a secondary description of
syntax, is $1-\max_{\ell}|S(e)\cap S(\ell)|/|S(e)\cup S(\ell)|$ with $S(e)$ the
set of structural hashes of $e$'s operator and terminal subtrees; it cannot
tell a re-spelled classic factor from new structure. The library has 18
entries covering momentum, reversal, volatility, liquidity and volume-price
interaction. Three are translations of published formulas of
\citet{kakushadze2016alphas} (Alphas \#2, \#6 and \#13); four more are
the project's own variants in that style.

\section{Experimental Setup}
\label{sec:setup}

\subsection{Synthetic planted-alpha benchmark}
\label{sec:synthetic}

\paragraph{Purpose.} The benchmark checks that the harness works end to end
(evidence level~2 of the suite's governance ladder). It measures whether a
search method can recover signals that are \emph{known} to exist, and how
often the protocol selects anything on \emph{null} markets where nothing is
planted, under the same budget and protocol as the real-data study. It says
nothing about real markets.

\paragraph{Data-generating process.} For asset $i$ on day $t$, the return is
\[
r_{t,i}=\beta_i m_t+\sigma_i\varepsilon_{t,i}
      +s\sum_{k=1}^{K}w_k\,z_{k;t-\ell-1,i},\qquad
z_{k;t}=\mathrm{cs\_zscore}\big(g_k(\text{bars up to }t)\big),
\]
where:
\begin{itemize}
  \item each $g_k$ is a planted expression in the factor language;
  \item $s=\mathrm{SNR}\cdot\overline{\sigma}$, with $\overline{\sigma}$ the
    average idiosyncratic volatility;
  \item the weights satisfy $\sum_k w_k^2=1$.
\end{itemize}
The planted component of $r_t$ uses signal values from $t-\ell-1$. A signal
observed at close $t$ therefore earns its return exactly after the execution
lag used by the evaluator.

The remaining parameters are:
\begin{itemize}
  \item $\beta_i\sim U(0.5,1.5)$, $\sigma_i\sim U(0.012,0.030)$,
    $m_t\sim N(0.0003,0.010^2)$ and $\varepsilon_{t,i}\sim N(0,1)$;
  \item returns are clipped at $\pm0.5$;
  \item log volume follows an AR(1) process with persistence $0.7$ and
    innovation scale $0.35$, of which a share of $0.3$ is common across
    assets. Volume is independent of returns;
  \item open, high and low are noise around the close path.
\end{itemize}
The planted signals depend on prices that themselves contain the planted
component. The price path is therefore solved as a fixed point: returns are
rebuilt and every $z_k$ is re-evaluated \emph{with the package's own
evaluator} until the maximum absolute change is at most $10^{-12}$. On the
planted markets this took
\SynIterationsLow{} iterations at the low SNR level and \SynIterationsHigh{}
at the high one, with a maximum residual of \SynMaxResidual{}. The bars have finality \code{confirmed} and pass the
same panel validation as real data.

\paragraph{Planted signals.} There are $K=4$ signals with equal weights, in
three families:
\begin{itemize}
  \item \emph{textbook}: \code{reversal\_5}, \code{-ts\_mean(returns, 5)},
    and \code{abnormal\_volume\_20}, \code{volume / ts\_mean(volume, 20)};
  \item \emph{library family}: \code{volume\_return\_corr\_10},
    \code{-ts\_corr(returns, delta(log(volume), 1), 10)}, the volume--return
    correlation idea behind Alpha\#2 of \citet{kakushadze2016alphas}, whose
    formula is in the reference library;
  \item \emph{drawn}: \code{\SynDrawnName}, \SynDrawnExpression{}, drawn at
    random from the typed grammar by a procedure fixed before the draw. The
    first candidate meeting four criteria was accepted: depth at least 3,
    complexity at most 8 and a time-series operator; well defined on a
    reference null market; mean rank correlation below $0.3$ in absolute
    value with every library entry and every other planted signal; and a
    convergent fixed point at both SNR levels. The draw record, with every
    rejected candidate and its reasons, is committed with the results.
\end{itemize}
The textbook signals favour any proposer with prior knowledge of classic
factors, such as an LLM, and so does the library-family signal to a lesser
degree. Only the drawn signal is unknown to every proposer in advance; it
comes from the random baseline's own sampling distribution, which if anything
favours that baseline. Being drawn at random does not prove that the
expression is absent from the literature; the correlation criterion only
shows that it is far, in values, from the reference library. Measured on the
formation window of the planted markets, behavioural novelty against the
library was \SynNoveltyReversal{} and \SynNoveltyVolume{} for the textbook
signals, \SynNoveltyLibraryFamily{} for the library-family signal and
\SynNoveltyDrawn{} for the drawn one. Every planted signal (up to its sign,
which the protocol chooses) has a non-zero sampling probability under the
random-grammar baseline, and its components are on genetic programming's
menus.

\paragraph{Grid and protocol.} The grid and protocol settings are:
\begin{itemize}
  \item \SynSymbols{} fictional symbols and \SynSessions{} sessions;
  \item SNR $\in\{\SynSnrLow,\SynSnrHigh\}$ and market seeds $\{0,1,2\}$,
    plus \SynNullSeeds{} null markets (SNR $0$, seeds $0$ to $9$); the
    proposer seed is the market seed plus $10{,}000$;
  \item a budget of \SynBudget{} trials and batches of \SynBatch{};
  \item $\ell=h=1$, fractions $0.6/0.2/0.2$ and an embargo of 2 sessions;
  \item a formation screen (BH at $q=\SynFdrAlpha$ over behavioural
    classes, $m$ = budget minus behavioural duplicates), BH confirmation at
    $q=\SynConfirmAlpha$ on validation, top-$k=\SynTopK$ and decorrelation at
    $|\rho|<\SynMaxCorr$;
  \item exactly one reveal per run.
\end{itemize}
The realized dispersion ratio of the planted component to idiosyncratic
noise was \SynRealizedSnrLow{} on average at the low SNR level and
\SynRealizedSnrHigh{} at the high one.

\paragraph{Benchmark metrics.}
\begin{itemize}
  \item \emph{Recovery}: planted signal $k$ counts as recovered if some
    selected factor's oriented test-window values have mean cross-sectional
    Spearman correlation $\bar\rho\ge\SynRecoveryThreshold$ with $z_k$. The
    rule is signed: an oriented factor that bets against $z_k$ recovers
    nothing.
  \item \emph{Recall}: $|\bar\rho|\ge\SynRecoveryThreshold$ for any
    evaluated trial on the formation window, using every
    \SynRecallStride{}th formation date. Recall asks whether the proposer
    ever found the signal, selected or not.
  \item \emph{True FDP}: the share of selected factors that are false by
    ground truth, i.e.\ whose oriented rank correlation with the planted
    composite (proportional to the true expected return) on the test window
    is below \SynTrueDiscoveryThreshold{}. On a null market every selected
    factor is false.
  \item \emph{Test non-significance rate}: the share of selected factors
    whose one-sided test IC is not significant under BH at \SynTestAlpha{}
    within the selected set. It measures the power of the test window, not
    falsity.
  \item Composite and mean per-factor test IC and ICIR, behavioural and
    structural novelty, the selection funnel (behavioural duplicates,
    degenerate trials, screen survivors, confirmed candidates) and, on null
    markets, the share of runs that selected anything.
\end{itemize}
As an oracle reference, we score the planted signals themselves and their
composite on the same test window; the oracle's \emph{power} is the share of
planted signals that pass BH at \SynTestAlpha{} on the test window.

\subsection{Real-market study (pre-registered; not yet run)}
\label{sec:real}

The real-market design is fixed in the project's research plan before any
test-window access. We summarize it here.

\paragraph{Data.} US equity daily bars from MarketData, read only through the
suite's gateway with finality \code{confirmed}.
\todo{Document the provider's split and dividend adjustment convention and
the data-quality report (coverage, stale prices, extreme returns).}

\paragraph{Universes.}
\begin{itemize}
  \item U1: liquid large capitalizations, selected by median formation-window
    dollar volume.
  \item U2: a broader liquid universe.
\end{itemize}
Membership is frozen at the end of the formation window.
\todo{Final sizes and the share of names with truncated histories (a
survivorship diagnostic).}

\paragraph{Windows.} The data are split chronologically into formation,
validation and sealed test windows with the embargo rule of
Section~\ref{sec:metrics}. The test window must extend beyond the documented
training-data cutoff of every model evaluated. A prospective window of
sessions after the commitment date follows.
\todo{Exact dates, fixed before the first reveal.}

\paragraph{Arms and budgets.} The arms are:
\begin{itemize}
  \item A: random grammar;
  \item B: genetic programming;
  \item C: the classic library passed through the same selection;
  \item D: LLM-guided search;
  \item optionally E: an interactive-mining prompting
    style~\citep{wang2023alphagpt} re-implemented in the same language and
    budget;
  \item F: compute-matched A and B, whose budget matches the LLM arm's
    measured cost (at least $10^4$ trials) under the same trial accounting.
\end{itemize}
The primary budget is $B=200$ trials, with $R=5$ replicates per arm and
universe. Equal trial budgets are not equal compute, and GP and
reinforcement-learning alpha miners~\citep{yu2023alphagen} typically use far
more evaluations, which is why arm F exists. Every planned test reveal is
listed in advance and counted in the study registry.

\paragraph{Hypotheses and decision rules.}
\begin{itemize}
  \item \textbf{H1.} The LLM beats both baselines on the sealed test. Two
    estimands are tested separately: the mean daily difference in pooled
    composite test IC (market sampling; a stationary block bootstrap over
    test dates, in the spirit of \citet{hansen2005spa}) and the difference in
    replicate-level composite ICIR (search-procedure variability; an
    unpaired exact Mann--Whitney test, since LLM and baseline replicates share
    nothing to pair on). Holm correction over the two comparisons; support
    also requires a pre-set minimum median improvement.
  \item \textbf{H2 (synthetic, seeds 100--119).} The LLM recovers the drawn
    signal more often. The test is an exact McNemar test paired by market
    seed, with Holm correction. The seeds exclude the harness seeds already
    seen.
  \item \textbf{H3.} The LLM's edge is not explained by period memorization.
    The decision uses only the difference-in-differences $D$ of pre- versus
    post-cutoff test IC between the LLM and the pre-specified GP baseline:
    memorization is flagged if a bootstrap interval for $D$ lies below zero,
    and absence of memorization is declared only by an equivalence test with
    a margin fixed in advance. It is declared underpowered with fewer than
    250 post-cutoff sessions.
  \item \textbf{H4.} Formation feedback improves search. It is tested on
    synthetic recall and on validation-window ICIR only.
\end{itemize}
Every test is one-sided at level $0.05$; the equivalence test of H3 is a pair of one-sided tests.

\paragraph{Ablations.} Ablations use validation data only. They vary:
\begin{itemize}
  \item feedback (on or off);
  \item the rationale requirement (on or off);
  \item anonymization (on or off, as a contamination probe in a pre-cutoff
    window only);
  \item reasoning effort;
  \item budget.
\end{itemize}

\subsection{Implementation and reproducibility}

\paragraph{Software.} Everything is implemented in Python. The recorded
software stack is \SynSoftware{}. The offline test suite checks:
\begin{itemize}
  \item absence of look-ahead for every operator;
  \item the information barrier;
  \item ledger tamper detection;
  \item the seal rules;
  \item the LLM backend's request shape and error handling, with a fake
    client.
\end{itemize}

\paragraph{Runs and artifacts.} The synthetic benchmark (\SynRuns{} planted
and \SynNullRuns{} null runs) took \SynRuntime{}~s of wall-clock time. It is
deterministic: ledger timestamps are disabled, and rerunning four of its runs
under a different hash seed reproduced their ledger heads and every scored
field. Each run stores:
\begin{itemize}
  \item its configuration and the hash-chained ledger;
  \item the selection and the reveal;
  \item a provenance file with the data hash, template hashes, served model
    identifiers, software versions and the ledger head.
\end{itemize}
LLM runs additionally record every request and its outcome, failures
included, for exact offline replay.

\section{Results}
\label{sec:results}

% Every number in this section comes from tables/*.tex, which
% scripts/render_paper_tables.py generates from
% results/synthetic_benchmark/summary.json. Do not type numbers by hand.
% LLM and real-data subsections stay as \todo placeholders until the
% pre-registered runs exist.

\subsection{Harness validation on synthetic data}
\label{sec:results-synthetic}

\noindent\fbox{\parbox{\dimexpr\linewidth-2\fboxsep-2\fboxrule}{\small
\textbf{Synthetic data --- not evidence about real markets.} This experiment
validates the search harness on simulated prices with known planted signals
and on null markets with none. Only the two baselines were run; the LLM arm
was not run.}}
\medskip

Table~\ref{tab:synthetic-headline} reports the headline metrics of the
\SynRuns{} planted runs (2 proposers $\times$ 2 SNR levels $\times$
\SynSeeds{} market seeds), Table~\ref{tab:synthetic-recovery} breaks recovery
down by planted signal, Table~\ref{tab:synthetic-null} reports the
\SynNullRuns{} null runs, and Table~\ref{tab:synthetic-oracle} gives the
planted signals' own test-window performance as a ceiling. Every run used
its full budget of \SynBudget{} trials (\SynIncompleteRuns{} incomplete
runs).

\begin{table}[t]
\centering
\caption{Synthetic planted-alpha benchmark, planted markets (mean $\pm$ sd
across complete runs; \SynBudget{} trials per run). Recovery, recall and the
number selected average all runs; the remaining columns describe selected
factors and average only the runs that selected something, so a value
without $\pm$ comes from one run. Composite = equal-weight z-scored
combination of the selected factors on the \SynTestIcDates{}-date sealed
test window. Synthetic data only.}
\label{tab:synthetic-headline}
\resizebox{\textwidth}{!}{% Generated by scripts/render_paper_tables.py from results/synthetic_benchmark/summary.json.
% Do not edit by hand. SYNTHETIC DATA - NOT EVIDENCE ABOUT REAL MARKETS.
\begin{tabular}{llcccccccccc}
\toprule
SNR & Arm & Complete & Recovery & Recall & Selected & True FDP & Test non-sig. & Composite test IC & Composite test ICIR & Behav. novelty & Struct. novelty \\
\midrule
0.05 & evolutionary & 3/3 & 0.08 $\pm$ 0.14 & 0.42 $\pm$ 0.14 & 1.3 $\pm$ 2.3 & 0.00 & 0.50 & 0.0224 & 0.162 & 0.38 & 0.88 \\
0.05 & random & 3/3 & 0.00 $\pm$ 0.00 & 0.42 $\pm$ 0.14 & 0.3 $\pm$ 0.6 & 0.00 & 0.00 & 0.0186 & 0.131 & 0.59 & 0.86 \\
0.15 & evolutionary & 3/3 & 0.33 $\pm$ 0.14 & 0.50 $\pm$ 0.00 & 5.0 $\pm$ 0.0 & 0.00 $\pm$ 0.00 & 0.00 $\pm$ 0.00 & 0.0762 $\pm$ 0.0126 & 0.581 $\pm$ 0.074 & 0.27 $\pm$ 0.07 & 0.85 $\pm$ 0.02 \\
0.15 & random & 3/3 & 0.33 $\pm$ 0.14 & 0.42 $\pm$ 0.14 & 5.0 $\pm$ 0.0 & 0.00 $\pm$ 0.00 & 0.13 $\pm$ 0.12 & 0.0677 $\pm$ 0.0040 & 0.540 $\pm$ 0.050 & 0.35 $\pm$ 0.08 & 0.82 $\pm$ 0.05 \\
\bottomrule
\end{tabular}
}
\end{table}

\begin{table}[t]
\centering
\caption{Recovery of each planted signal: runs in which a \emph{selected}
factor matched it on the test window ($\bar\rho\ge\SynRecoveryThreshold$),
and runs in which \emph{any evaluated trial} matched it on the formation
window ($|\bar\rho|\ge\SynRecoveryThreshold$; ``found'', i.e.\ recall, which
is about as frequent on the null markets). Synthetic data only.}
\label{tab:synthetic-recovery}
\resizebox{\textwidth}{!}{% Generated by scripts/render_paper_tables.py from results/synthetic_benchmark/summary.json.
% Do not edit by hand. SYNTHETIC DATA - NOT EVIDENCE ABOUT REAL MARKETS.
\begin{tabular}{llcccc}
\toprule
SNR & Arm & \texttt{reversal\_5} & \texttt{abnormal\_volume\_20} & \texttt{volume\_return\_corr\_10} & \texttt{drawn\_40} \\
\midrule
0.05 & evolutionary & 0/3 selected, 3/3 found & 1/3 selected, 2/3 found & 0/3 selected, 0/3 found & 0/3 selected, 0/3 found \\
0.05 & random & 0/3 selected, 3/3 found & 0/3 selected, 2/3 found & 0/3 selected, 0/3 found & 0/3 selected, 0/3 found \\
0.15 & evolutionary & 2/3 selected, 3/3 found & 2/3 selected, 3/3 found & 0/3 selected, 0/3 found & 0/3 selected, 0/3 found \\
0.15 & random & 2/3 selected, 3/3 found & 2/3 selected, 2/3 found & 0/3 selected, 0/3 found & 0/3 selected, 0/3 found \\
\bottomrule
\end{tabular}
}
\end{table}

\begin{table}[t]
\centering
\caption{Null markets (SNR $0$, \SynNullSeeds{} seeds): runs that selected
at least one factor, and the selection funnel (mean $\pm$ sd across runs).
Every factor selected here would be a false discovery. Synthetic data only.}
\label{tab:synthetic-null}
\resizebox{\textwidth}{!}{% Generated by scripts/render_paper_tables.py from results/synthetic_benchmark/summary.json.
% Do not edit by hand. SYNTHETIC DATA - NOT EVIDENCE ABOUT REAL MARKETS.
\begin{tabular}{lcccccc}
\toprule
Arm & Complete & Runs selecting $\geq 1$ & Selected & Screen survivors & Confirmed & Composite test $t$ \\
\midrule
evolutionary & 10/10 & 0/10 & 0.0 $\pm$ 0.0 & 4.5 $\pm$ 10.9 & 0.0 $\pm$ 0.0 & n/a \\
random & 10/10 & 0/10 & 0.0 $\pm$ 0.0 & 0.2 $\pm$ 0.6 & 0.0 $\pm$ 0.0 & n/a \\
\bottomrule
\end{tabular}
}
\end{table}

\begin{table}[t]
\centering
\caption{Oracle reference: test-window mean IC (Newey--West $t$) of the planted
signals and of their equal-weight composite, per planted market, and the
share of planted signals that pass BH at \SynTestAlpha{} on the test window
(power). Synthetic data only.}
\label{tab:synthetic-oracle}
\resizebox{\textwidth}{!}{% Generated by scripts/render_paper_tables.py from results/synthetic_benchmark/summary.json.
% Do not edit by hand. SYNTHETIC DATA - NOT EVIDENCE ABOUT REAL MARKETS.
\begin{tabular}{llcccccc}
\toprule
SNR & Seed & \texttt{reversal\_5} & \texttt{abnormal\_volume\_20} & \texttt{volume\_return\_corr\_10} & \texttt{drawn\_40} & Planted composite & Power \\
\midrule
0.05 & 0 & 0.0180 (1.90) & 0.0407 (4.23) & 0.0177 (1.64) & 0.0238 (2.25) & 0.0554 (4.72) & 0.75 \\
0.05 & 1 & 0.0274 (2.59) & 0.0300 (2.56) & 0.0244 (2.28) & 0.0264 (2.54) & 0.0522 (4.65) & 1.00 \\
0.05 & 2 & 0.0160 (1.69) & 0.0424 (3.82) & 0.0258 (2.64) & 0.0308 (3.09) & 0.0554 (4.90) & 1.00 \\
0.15 & 0 & 0.0537 (5.71) & 0.0830 (8.80) & 0.0562 (5.52) & 0.0416 (3.88) & 0.1333 (11.94) & 1.00 \\
0.15 & 1 & 0.0652 (6.39) & 0.0700 (5.86) & 0.0625 (5.92) & 0.0568 (5.12) & 0.1370 (12.16) & 1.00 \\
0.15 & 2 & 0.0506 (5.19) & 0.0870 (7.39) & 0.0701 (7.44) & 0.0681 (7.36) & 0.1445 (12.95) & 1.00 \\
\bottomrule
\end{tabular}
}
\end{table}

We draw six observations from the tables. With \SynSeeds{} seeds per planted
cell and \SynNullSeeds{} null seeds, we read them as descriptions of the
harness and the baselines, not as statistical claims.

\paragraph{The confirmation step, not the formation screen, controls false
selections.} On the null markets neither arm selected a factor
(\SynNullSelecting{evolutionary} runs for GP and \SynNullSelecting{random}
for random search). The formation screen alone was not safe for the
feedback-driven arm: it passed \SynNullMeanSd{evolutionary}{n_screen_survivors}
expressions per null run for GP against
\SynNullMeanSd{random}{n_screen_survivors} for random search, and none of
them was confirmed on validation. GP chooses its next expressions after
seeing formation results, so its formation p-values are biased toward
significance. This is why the protocol treats the formation screen as a
filter and places the claim-carrying multiplicity step on data no proposer
has seen.

\paragraph{The protocol has little power at the low signal-to-noise level.}
At SNR \SynSnrLow{}, GP selected \SynMeanSd{0.05}{evolutionary}{n_selected}
factors per run and random search \SynMeanSd{0.05}{random}{n_selected}; the
share of runs that selected anything was
\SynMean{0.05}{evolutionary}{any_selected} and
\SynMean{0.05}{random}{any_selected}. The confirmation step confirmed
\SynMean{0.05}{evolutionary}{n_confirmed} and
\SynMean{0.05}{random}{n_confirmed} candidates on average, out of
\SynMean{0.05}{evolutionary}{n_screen_survivors} and
\SynMean{0.05}{random}{n_screen_survivors} screen survivors. The planted
composite itself is clearly detectable on the test window at this level
(Table~\ref{tab:synthetic-oracle}), but the individual planted signals are
weaker, and the validation window holds only a fifth of the sessions. The
real-data design therefore uses multi-year windows.

\paragraph{At the higher level, selections are true discoveries but mostly
known ones.} At SNR \SynSnrHigh{} GP selected
\SynMeanSd{0.15}{evolutionary}{n_selected} factors per run and random search
\SynMeanSd{0.15}{random}{n_selected}, and the true FDP was \SynMeanSd{0.15}{evolutionary}{fdp_true} for GP
and \SynMeanSd{0.15}{random}{fdp_true} for random search. The test
non-significance rate, \SynMeanSd{0.15}{evolutionary}{test_nonsignificant_rate}
and \SynMeanSd{0.15}{random}{test_nonsignificant_rate}, therefore measures
power, not error. The composite test IC was
\SynMeanSd{0.15}{evolutionary}{composite_test_ic} (GP) and
\SynMeanSd{0.15}{random}{composite_test_ic} (random), against
\SynOracleHighMin{}--\SynOracleHighMax{} for the planted composite; the
difference between the arms is within one seed standard deviation. Only the
textbook signals were recovered: neither arm selected the library-family
signal (\SynLibraryFamilySelected{} runs) or the drawn signal
(\SynDrawnSelected{} runs), and no evaluated trial matched either
(\SynLibraryFamilyFound{} and \SynDrawnFound{} runs), although both are in
the baselines' search space. At this budget, the part of the benchmark that
prior knowledge cannot solve is untouched by undirected search.

\paragraph{Syntax-based novelty would have misled.} The factors selected at
SNR \SynSnrHigh{} have structural novelty
\SynMeanSd{0.15}{evolutionary}{structural_novelty_mean} (GP) and
\SynMeanSd{0.15}{random}{structural_novelty_mean} (random), which would read
as new structure. Their behavioural novelty is
\SynMeanSd{0.15}{evolutionary}{behavioural_novelty_mean} and
\SynMeanSd{0.15}{random}{behavioural_novelty_mean}: their values are
substantially correlated with those of library factors (behavioural novelty
is one minus the largest such correlation), although only a few are exact
re-spellings. Rediscovery must be measured on values, which is how the LLM
study will report it.

\paragraph{Recall does not measure detection.} A planted expression counts as
``found'' when any evaluated trial matches it on the formation window. On the
null markets, where the same expressions carry no return, recall was
\SynNullMeanSd{evolutionary}{recall_rate} for GP and
\SynNullMeanSd{random}{recall_rate} for random search, against
\SynMeanSd{0.15}{evolutionary}{recall_rate} and
\SynMeanSd{0.15}{random}{recall_rate} on the planted markets at SNR
\SynSnrHigh{}. Recall therefore records that a proposer writes such
expressions (close variants of the short-reversal and abnormal-volume
signals are easy to reach in the grammar), not that it detected a signal; H4 uses recall only as a contrast
with null markets.

\paragraph{Syntactic deduplication undercounts repeated hypotheses.} Genetic
programming produced \SynMeanSd{0.15}{evolutionary}{n_behaviour_duplicates}
evaluated trials per run at SNR \SynSnrHigh{} whose per-date ranks repeated
an earlier trial's, and random search
\SynMeanSd{0.15}{random}{n_behaviour_duplicates}. Counted as separate
hypotheses, such clusters would inflate the formation screen's discoveries,
since a step-up procedure rewards clusters of equal p-values.

\subsection{LLM-guided search on the synthetic benchmark}
\label{sec:results-llm-synthetic}

\todo{Pending: no LLM experiment has been run. After the pre-registered run
(seeds 100--119 $\times$ 2 SNR levels plus null markets, responses recorded
for replay), report the following: the LLM arm in
Tables~\ref{tab:synthetic-headline} to~\ref{tab:synthetic-null}; the H2
McNemar tests on the drawn signal; validity, duplicate and degenerate rates;
recall as a function of trials and of cost, against the compute-matched
arm; served model identifiers; and token usage.}

\subsection{Real-market comparison (RQ1, RQ2)}
\label{sec:results-real}

\todo{Pending: no real-market evaluation has been run. After the
pre-registered reveals, report the following per arm, universe and
replicate: trial accounting and the selection funnel; sealed-test IC, ICIR
and Newey--West $t$ of the composite and of each factor; the test
non-significance rate; long-short net Sharpe with cost sensitivity (0, 5, 10
and 20 bps); turnover; the registry's exploration and reveal counts; and the
H1 tests with their decision.}

\subsection{Memorization and look-ahead (RQ3)}
\label{sec:results-contamination}

\todo{Pending: report the pre/post knowledge-cutoff difference-in-differences
against GP (H3) with the cutoff date's documentation source, the
behavioural-novelty-stratified LLM advantage, and the anonymization-probe
results on the pre-cutoff validation window.}

\subsection{Ablations (RQ4)}
\label{sec:results-ablations}

\todo{Pending: report the feedback on/off test (H4), the rationale on/off and
effort ablations, and the budget-scaling study, on synthetic data and the
validation window only.}

\section{Discussion}
\label{sec:discussion}

\paragraph{What the protocol buys.} The design choices trade statistical
power for credibility.
\begin{itemize}
  \item \emph{Counting every trial.} Invalid proposals, evaluation failures
    and degenerate trials enter the family size, and so does unused budget.
    This penalizes proposers that generate many bad candidates or stop
    early, which is exactly the behaviour that inflates apparent discovery
    rates in automated search.
  \item \emph{Counting each hypothesis once.} Trials with identical per-date
    ranks are one hypothesis. Syntactic deduplication alone would let a
    proposer inflate its discoveries by re-spelling a good expression.
  \item \emph{Confirming on unseen data.} A proposer that learns from
    formation feedback makes formation p-values optimistic; the null markets
    of Section~\ref{sec:results-synthetic} show it for genetic programming.
    Only the validation step, on data no proposer has seen, carries
    false-discovery control.
  \item \emph{The information barrier.} The harness passes proposers
    formation statistics only, and on store data later windows are not
    loaded until their phase. For an LLM, whose only input is the rendered
    prompt, this means the model cannot condition on validation or test
    outcomes of \emph{this} study. It does not remove what the model may
    have learned about the period in pre-training; that is what the
    contamination diagnostics address.
  \item \emph{The commit-then-reveal hold-out.} The ledger and the study
    registry give tamper evidence for one's own records: every commitment,
    reveal and refusal is chained, and repeated reveals are counted. They do
    not by themselves prove to a third party that no earlier reveal or
    discarded run exists, because the same process writes them. That proof
    needs the commitment hash published outside the repository, for example
    in a pushed signed tag or a public timestamp, before the separate reveal
    step, which the command-line interface enforces for real data.
\end{itemize}
The controls are designed to apply equally to every proposer, but they are
not neutral in effect: equal trial budgets are not equal compute, which is
why the design includes a compute-matched baseline arm. They make results
harder to obtain and easier to believe.

\paragraph{Reading the synthetic results.} The harness results are a
statement about the \emph{baselines} and the protocol. At a budget of
\SynBudget{} trials, undirected random search and feedback-driven genetic
programming both reach the textbook planted signals at the higher
signal-to-noise level, and both miss the library-family and the drawn
signals, which are in their search space. The selected factors are true
discoveries by the benchmark's ground truth, but mostly substantially
correlated with library factors. The benchmark thus has room to show an advantage from
better proposals, and a test of whether an advantage is more than prior
knowledge.

Such an advantage would only be informative under three conditions:
\begin{enumerate}
  \item it appears on the drawn signal, not only on the textbook and
    library-family ones that an LLM may know from the literature;
  \item it survives the same multiple-testing accounting and confirmation;
  \item it replicates across the pre-registered seeds (H2) and holds against
    compute-matched baselines.
\end{enumerate}

\paragraph{What would count as evidence for LLM-guided search.} By the
pre-registered rules, a positive answer to RQ1 requires three things:
\begin{itemize}
  \item a sealed-test improvement over both baselines at an equal budget, on
    both the market-sampling and the search-variability estimands (H1);
  \item an equivalence result showing that the improvement does not shrink
    after the model's training cutoff relative to genetic programming (H3);
  \item ideally, confirmation on a prospective window that did not exist
    when the factor set was committed.
\end{itemize}
An improvement confined to factors with low behavioural novelty would count
as rediscovery, not discovery. \todo{Discuss the actual outcome of H1 to H4
once the pre-registered runs exist.}

\paragraph{Rationales are claims, not evidence.} The LLM proposer attaches an
economic rationale and a mechanism label to every expression. These texts
are recorded but play no role in selection. Fluent but unsupported
explanations are a known failure mode of language
models~\citep{ji2023hallucination}. Whether rationales carry information
about out-of-sample survival is itself an empirical question for the
rationale ablation. \todo{Report whether rationale presence or content
predicts test survival.}

\paragraph{Scope and use.} This is research software. It places no orders
and makes no claim that any discovered factor would remain profitable after
costs, capacity limits or publication. Market prices in the real-data study
come only from a commercial data provider (MarketData) through the suite's
gateway, and no vendor data are redistributed.

\section{Limitations}
\label{sec:limitations}

\paragraph{Evidence status.} The only completed experiment is the synthetic
harness validation, with \SynSeeds{} seeds per planted cell, \SynNullSeeds{}
null markets and two baselines. No LLM experiment and no real-market
evaluation has been run. All claims about LLM-guided search in this draft are
hypotheses. The benchmark's planted signals and hypothesis H2 were revised
after the first baseline harness runs had been seen (Section~\ref{sec:setup}
and the research plan disclose this); the confirmatory seeds exclude the
seeds already seen.

\paragraph{Synthetic benchmark design.} The authors chose the following, and
rankings on this benchmark need not transfer to real markets:
\begin{itemize}
  \item the data-generating process, the textbook and library-family planted
    signals, the drawing procedure of the drawn signal, the SNR levels and
    the genetic-programming settings;
  \item recovery depends on a correlation threshold of
    \SynRecoveryThreshold{} and the true FDP on a threshold of
    \SynTrueDiscoveryThreshold{};
  \item recall uses every \SynRecallStride{}th formation date for speed.
\end{itemize}
Returns and volumes are independent apart from the planted terms. The
volume process is not calibrated to any market, and there are no
listings, delistings or missing data. The validation and test windows are
short, so power at the lower signal-to-noise level is low.

\paragraph{Statistical simplifications.}
\begin{itemize}
  \item p-values refer the Newey--West $t$-statistic to a $t$ distribution
    with $n-1$ degrees of freedom, a convention rather than an exact
    small-sample result.
  \item The deflated Sharpe ratio is a diagnostic that treats trials as
    independent, with $2N$ trials for the two-sided selection on $|\ICIR|$.
  \item Benjamini--Hochberg control is valid under
    independence~\citep{benjamini1995fdr} or positive regression
    dependence~\citep{benjamini2001fdr}. Behavioural deduplication merges
    only trials with identical per-date ranks; highly correlated variants
    remain separate hypotheses, and for a step-up procedure such clusters
    make the screen less strict. The Benjamini--Yekutieli variant is
    therefore reported as well.
  \item Formation p-values of feedback-driven proposers are not valid for
    false-discovery control; the protocol relies on the validation
    confirmation for its error control.
  \item Bootstrap data-snooping tests~\citep{white2000reality,hansen2005spa,romano2005stepwise},
    and the block bootstrap that the pre-registered H1 and H3 use, are not
    yet implemented.
\end{itemize}

\paragraph{Protocol gaps and trust assumptions.}
\begin{itemize}
  \item The trial ledger and the study registry are self-attested. They
    detect accidental or naive edits and make repeated reveals countable;
    they rely on commitment hashes being published outside the repository
    before a reveal.
  \item The information barrier is an interface, not a sandbox. In-process
    code could read anything in memory, and for synthetic data the full
    panel is in memory.
  \item The classic-library baseline, the compute-matched baseline and the
    ablation wrappers are not implemented or not yet run.
\end{itemize}

\paragraph{Contamination controls are partial.}
\begin{itemize}
  \item Anonymization removes dates, tickers and sample sizes from prompts,
    and a render-time check guards the harness-written text. The check is
    heuristic: feedback statistics are not checked for year-like numbers,
    distributional fingerprints of a period cannot be detected, and generic
    knowledge of documented anomalies cannot be removed.
  \item The knowledge-cutoff split depends on the provider's documentation
    of the training-data cutoff and on enough post-cutoff sessions for
    power.
  \item Only a prospective window is out of sample for the model by
    construction.
\end{itemize}

\paragraph{Real-data threats.}
\begin{itemize}
  \item Without point-in-time index membership or delisted histories, the
    universes may be affected by survivorship bias.
  \item Close-to-close returns include or omit dividends depending on the
    provider's adjustments.
  \item \code{vwap\_proxy} is a typical price, not a traded VWAP.
  \item The long-short evaluation uses a flat cost of 10 basis points per
    unit of turnover. It ignores borrow costs, market impact and capacity,
    and it trades at vendor closes rather than official auction prints.
\end{itemize}

\paragraph{LLM-specific issues.}
\begin{itemize}
  \item Current models accept no sampling parameters, so replicates vary in
    ways we record but do not control.
  \item Recorded outcomes make each run exactly replayable, and record files
    are write-only, so a new live run is always a new sample.
  \item Served model identifiers are logged, and server-side fallback is
    disabled, because it could otherwise change the model under study
    silently.
\end{itemize}


\bibliographystyle{plainnat}
\bibliography{references}

\end{document}
